\documentclass[10pt,a4paper]{article}
\usepackage[margin=2cm]{geometry}
\usepackage{booktabs,array,longtable}
\usepackage[table]{xcolor}
\usepackage{titlesec}
\usepackage{enumitem}
\usepackage[colorlinks,allcolors=blue!55!black]{hyperref}
\usepackage{parskip}
\definecolor{hd}{RGB}{18,48,88}
\definecolor{bad}{RGB}{150,20,20}
\definecolor{good}{RGB}{20,110,45}
\titleformat{\section}{\large\bfseries\color{hd}}{\thesection}{0.6em}{}
\titleformat{\subsection}{\normalsize\bfseries\color{hd}}{\thesubsection}{0.6em}{}
\newcommand{\B}[1]{\textcolor{bad}{\textbf{#1}}}
\newcommand{\G}[1]{\textcolor{good}{\textbf{#1}}}
\setlist{itemsep=1.5pt,topsep=2pt,parsep=0pt}
\renewcommand{\arraystretch}{1.18}
\title{\vspace{-1.4cm}\textbf{COMPASS Real-Data Specification}\\[1.5mm]
\large What the current corpus can and cannot support, and exactly what to collect}
\date{8 September 2026}\author{}
\begin{document}\maketitle\vspace{-9mm}\hrule\vspace{3mm}

\section{The one defect that matters}

The corpus is large. It is not the size that limits us.

Every real reconstruction number is produced by holding out one reference point and
predicting it from the others. Under that protocol the model is only ever asked to
interpolate between nearby samples of the same process, which is the regime where a
tuned interpolator is close to optimal by construction. The two things our method is
actually for, extrapolating into unmeasured regions and reasoning about geometry behind
walls, are invisible to it.

That is why classical wins on real data and learned methods lose. It is a property of
the measurement geometry, not of the models. \textbf{The fix is a densely sampled
reference floor that permits contiguous block hold-out.} Everything else in this document
is secondary to that.

Concrete evidence that the protocol is degenerate: in
\texttt{results/realdata\_recon.json} the \texttt{buffer\_0m} and \texttt{buffer\_1m}
results are \B{byte-identical}. The 1\,m exclusion buffer removed no test points at all,
because reference points are never within 1\,m of each other. At 2\,m the buffer finally
bites and every method degrades by about 27\%.

\section{What exists, precisely}

\subsection{Coverage is far thinner than the row count suggests}

\begin{center}\small
\begin{tabular}{@{}lrrrrrr@{}}
\toprule
Site / floor & RPs & Cells & Cells $\ge$8 RP & Extent (m) & Area (m$^2$) & RP spacing \\ \midrule
CMU-Q f1 & 84 & 25 & 13 & $41.1\times29.2$ & 1200 & 3.8\,m \\
\rowcolor{green!8}CMU-Q f2 & 84 & 35 & 15 & $39.7\times30.2$ & 1199 & 3.8\,m \\
CMU-Q f3 & 70 & 41 & 15 & $31.3\times30.1$ & 942 & 3.7\,m \\
EC Parking f0 & 20 & 18 & 9 & \B{$0.0\times0.0$} & \B{---} & \B{no coords} \\
EC Parking f1 & 34 & 27 & 15 & $48.6\times11.8$ & 573 & 4.1\,m \\
EC Parking f2 & 34 & 22 & 9 & $48.6\times11.8$ & 573 & 4.1\,m \\
\midrule
Ezdan Tower 1 & 2--5 per floor & 9--24 & \B{0 on every floor} & degenerate & --- & --- \\
Ezdan Tower 4 & 4 per floor & 4--17 & \B{0 on every floor} & --- & --- & --- \\
\bottomrule
\end{tabular}
\end{center}

\textbf{The two high-rises contribute nothing.} Ezdan Towers 1 and 4 hold 927{,}000
stationary rows across 55 floors and yield \B{zero} usable cells, because each floor has
only 2 to 5 reference points and a cell needs at least 8 to be reconstructable. All 76
viable cells come from CMU-Q (43) and EC Parking (33). Any statement about ``four sites,
61 floors'' is true of the raw capture and false of the evaluable data.

\subsection{What is present and what is absent}

Present per sample: \texttt{phoneName}, \texttt{rpNumber}, \texttt{scanNumber},
\texttt{timeStamp}, \texttt{transmitter\_id}, \texttt{transmitter\_rss},
\texttt{coordinates}, \texttt{floor}, plus floor plans at 0.032\,m per pixel.

Absent, in the order that hurts:

\begin{center}\small
\begin{tabular}{@{}p{4.2cm}p{12.6cm}@{}}
\toprule
Missing & Consequence today \\ \midrule
\B{Dense ground-truth field} & Cannot measure extrapolation. Forces leave-one-point-out,
which structurally favours interpolation. This single gap is why the real-data section
currently reads as a negative result. \\
\B{Ground-truth transmitter positions} & Only 17.9\% of cells pass the source-reliability
gate, and even those localise only to 0.26--9.19\,m by self-consistency. There is no
external reference, so no localisation error can be reported, and the occlusion field
cannot be computed on real data because it needs a transmitter to cast a ray from. \\
\B{Inertial signals} & No accelerometer, gyroscope, magnetometer or barometer anywhere.
The order argument therefore rests on a statistical control (shuffling) rather than on a
physical one (motion). \\
\B{Continuous trajectory position} & Mobile walks exist (1.33M rows) but carry
\texttt{rpNumber} $=-1$, so they have no position labels. The order and active-sensing
work uses them as sequences without knowing where they are. \\
Device calibration reference & Offsets are inferred, not measured against a controlled
capture. \\
Radio metadata & No band, PCI, or indoor-versus-macro flag, which is what would explain
the shared-identifier problem rather than merely asserting it. \\
Temporal repeats & No route walked on separate days, so drift is unquantified. \\
\bottomrule
\end{tabular}
\end{center}

\section{Target: one dense reference floor}

\subsection{Choose CMU-Q floor 2}

It has the most transmitters (35), the widest device spread in the corpus (13.5\,dB),
six phones sharing cells, a valid floor plan and coordinate file, and 600 scans per
reference point. It is also already the best-sampled floor, so the dense survey extends
rather than replaces existing work.

\subsection{Grid density and effort}

The floor is $39.7\times30.2$\,m, about 1199\,m$^2$ gross. Assuming 55--65\% is
accessible free space, roughly 700\,m$^2$ is walkable.

\begin{center}\small
\begin{tabular}{@{}lrrl@{}}
\toprule
Grid pitch & Points ($\approx$) & vs current 84 RPs & Assessment \\ \midrule
3.8\,m (current) & 84 & $1\times$ & Cannot support block hold-out. \\
2.0\,m & 175 & $2.1\times$ & Minimum viable. Blocks of 3--4\,m become possible. \\
\rowcolor{green!8}1.5\,m & 310 & $3.7\times$ & \G{Recommended.} Good block resolution, one to two days. \\
1.0\,m & 700 & $8.3\times$ & Ideal, but three to four days. Diminishing returns. \\
\bottomrule
\end{tabular}
\end{center}

At 1.5\,m, with a 45\,s dwell and 30\,s to move and mark, 310 points is about 6.5 hours
of collection. Two people with all seven phones on a shared rig finish in \G{two days}
including setup. That is the entire cost of turning the weakest section of the paper into
its strongest.

\subsection{Why this changes the result and not just the protocol}

With a dense field you can hold out \emph{contiguous regions} rather than scattered
points. Mask a 6\,m block, or a whole room, or everything beyond a given wall, and ask
each method to predict it. Interpolators have nothing to interpolate from inside the
mask, so they fall back to distance-weighted smoothing, while a geometry-conditioned
model can reason about what the wall does. That is the comparison our method was built
for and the one the current protocol cannot express.

It also, for the first time, makes the occlusion channel computable on real data, provided
the transmitter survey below happens.

\section{Collection specification}

\subsection{A. Dense grid survey (blocking)}

\begin{itemize}
\item CMU-Q floor 2, 1.5\,m grid across all accessible free space, roughly 310 points.
\item Mark the grid physically. Chalk plus a laser rangefinder from two known walls is
sufficient. Target positional error under 0.3\,m, which is a fifth of the pitch.
\item Dwell 45\,s per point. Log \emph{every} visible cell, not only strong ones, and
record non-detections explicitly. A cell that is not heard is information.
\item Carry all seven phones simultaneously on a rigid rig with fixed relative geometry,
so device comparisons at each point are controlled.
\item Record rig orientation (one of four cardinal headings) and repeat one in ten points
at a second heading to quantify body and orientation effects.
\end{itemize}

\subsection{B. Transmitter survey (unlocks the geometry claim)}

\begin{itemize}
\item Physically locate every accessible antenna, femtocell, booster and repeater on the
floor. Record position in the same coordinate frame as the grid, plus mounting height.
\item Record which cell identifiers each physical unit broadcasts. This is the
measurement that turns ``boosters share identifiers'' from an inference into a fact.
\item Twenty surveyed units is enough. This is the item that makes the ray-occlusion
field computable on real data, so it is second only to the grid itself.
\end{itemize}

\subsection{C. Trajectories with recoverable position}

\begin{itemize}
\item Walk routes that start, end and pass through marked grid points, pressing a marker
button at each, so position between fixes is interpolable rather than guessed.
\item Four route types: a corridor loop, a route crossing at least four walls, a route
into a deep non-line-of-sight pocket, and an open-space control.
\item Three phones carried together per walk, three times of day, three separate days.
\item Vary speed deliberately: slow, normal, brisk. Speed aliasing is a real confound in
the order argument and nobody has measured it.
\end{itemize}

\subsection{D. Inertial and context streams}

Logged continuously alongside every mobile session:

\begin{center}\small
\begin{tabular}{@{}llp{8.4cm}@{}}
\toprule
Stream & Minimum rate & Purpose \\ \midrule
Accelerometer & 50\,Hz & Step detection, dead reckoning between fixes \\
Gyroscope & 50\,Hz & Heading change, turn detection \\
Magnetometer & 25\,Hz & Absolute heading, indoor anomaly features \\
Barometer & 1\,Hz & Floor transitions \\
Orientation quaternion & 25\,Hz & Body and hand pose \\
Step counter & per event & Cheap odometry cross-check \\
Carry mode & per session & hand, pocket, bag \\
Monotonic + wall clock & per sample & Cross-device drift recovery \\
\bottomrule
\end{tabular}
\end{center}

\subsection{E. Device calibration capture}

All seven phones, same rig, same orientation, 2\,min dwell, at five points spanning
strong to weak signal. Repeat at the start and end of the campaign. This yields a
ground-truth offset \emph{as a function of signal level}, which matters because the
current model assumes a single additive scalar per device and the audit already flags
that real handsets may compress at strong signal.

\section{What each item buys}

\begin{center}\small
\begin{tabular}{@{}p{3.4cm}p{5.0cm}p{8.0cm}@{}}
\toprule
Collect & Unlocks & Claim it makes possible \\ \midrule
Dense grid & Block hold-out evaluation & ``The simulation ranking does not survive
transfer'' becomes measurable rather than asserted. Extrapolation becomes visible. \\
TX survey & Occlusion field on real data & ``Our geometric prior transfers'' becomes
testable. Currently untestable at any price. \\
Trajectory position & Order and active sensing on located walks & Turns the order result
from a sequence-imputation proxy into a spatial claim. \\
IMU & Motion-grounded order & Replaces the shuffle control with a physical one, and makes
the dataset multimodal, which is what a vision audience recognises. \\
Calibration capture & Level-dependent device model & Fixes the additive-scalar assumption
the audit flags as unrealistic. \\
Repeats & Temporal drift & Answers the generalisation question reviewers always ask. \\
\bottomrule
\end{tabular}
\end{center}

\section{Minimum viable versus ideal}

If time collapses, collect in this order and stop wherever you must. Each prefix is
independently publishable.

\begin{enumerate}
\item \textbf{Dense grid on one floor, plus transmitter survey.} This alone rescues the
real-data section and makes the geometry claim testable. Two to three days.
\item \textbf{Add located trajectories with IMU.} Makes the order and acquisition claims
spatial. Two more days.
\item \textbf{Add device calibration and temporal repeats.} Closes the two remaining
reviewer questions. One more day.
\item \textbf{Second dense floor} (CMU-Q f1 or f3) for a within-site generalisation axis.
\end{enumerate}

\textbf{A note on scope.} Do not spread thin across the two Ezdan towers. Fifty-five
floors there currently yield zero evaluable cells, and bringing even one of them to
viability needs the same dense survey as CMU-Q f2. One dense floor beats sixty sparse
ones for every claim in this paper.

\end{document}
